\subsection{Arbitrary fpqc descent and the main counterexample}

The following local algebra argument closes the descent gate isolated
in Proposition~\ref{A:hensel-remaining-gate}.  All base rings in this
subsection are commutative; neither the local algebra argument nor the
descent theorem imposes a noetherian hypothesis.

\begin{lemma}[A finite module containing the descent datum]
\label{A:fpqc-finite-module}
\status{Proved}
Let $A\to B$ be flat, let $b\in B$ satisfy a monic polynomial
$f\in A[T]$, and put $C=A[T]/(f)$.  Give $B$ the $C$-module
structure in which $T$ acts as multiplication by $b$.  Suppose that
$e\in B\otimes_A B$ satisfies
\[
 (b\otimes1-1\otimes b)e=0,\qquad \mu(e)=1,
\]
where $\mu$ is multiplication.  There is a $C$-module $M$, finitely
presented as an $A$-module, for which $B$ with this action is a direct
summand of $M\otimes_A B$ as a $C\otimes_A B$-module.
\end{lemma}
\begin{proof}
Write $e=\sum_{i=1}^r u_i\otimes v_i$.  Map $M_0=C^r$ to $B$
by sending its $i$th basis vector to $u_i$, and call the kernel $K$.
In $M_0\otimes_A B$ set $e_0=\sum_i e_i\otimes v_i$.
The image of $(T\otimes1-1\otimes b)e_0$ in $B\otimes_A B$
is zero.  Flatness of $B$ identifies the kernel of
$M_0\otimes_A B\to B\otimes_A B$ with $K\otimes_A B$.
Thus the displayed error is a finite sum $\sum_j k_j\otimes w_j$
with $k_j\in K$.  Let $K_0\subset K$ be the $C$-submodule
generated by these finitely many $k_j$, and put $M=M_0/K_0$.
The algebra $C$ is finite free over $A$, so $K_0$ is finitely
generated over $A$ and $M$ is finitely presented over $A$.

The map $g:M\to B$ induced by $M_0\to B$ is $C$-linear.
The image $e_M$ of $e_0$ satisfies $(T-b)e_M=0$ and
$g_B(e_M)=1$, where
\[
 g_B:M\otimes_A B\longrightarrow B,\qquad
       m\otimes z\longmapsto g(m)z.
\]
Consequently $z\mapsto e_Mz$ is a $C\otimes_A B$-linear
section of $g_B$, as required.
\end{proof}

\begin{proposition}[Indecomposable summands over a strictly henselian base]
\label{A:fpqc-indecomposable-basechange}
\status{Proved}
Let $(A,\mathfrak m,k)$ be strictly henselian local, let $M$ be a
finitely presented $A$-module with an endomorphism $T$, and let
$A\to(B,\mathfrak n,K)$ be a flat local homomorphism.
There is a finite direct-sum decomposition $M=\bigoplus_i M_i$
preserved by $T$ such that both
\[
 \operatorname{End}_{A[T]}(M_i)\quad\hbox{and}\quad
 \operatorname{End}_{B[T]}(M_i\otimes_A B)
\]
are local rings.  Zero summands are omitted.
\end{proposition}
\begin{proof}
We give the endomorphism-ring argument, including the points which
do not follow from noetherian module theory.  For any finite
$A$-module $N$ with an endomorphism $T$, put
\[
 E=\operatorname{End}_{A[T]}(N),\qquad
 J=\ker\bigl(E\longrightarrow\operatorname{End}_k(N/\mathfrak mN)\bigr).
\]
The image $R=E/J$ is a finite-dimensional $k$-algebra.
Moreover $J\subset\operatorname{Jac}(E)$.  Indeed, if $h\in J$
and $g\in E$, then $1-gh$ is the identity modulo $\mathfrak m$,
so its cokernel is zero by Nakayama's lemma.  A surjective
endomorphism $u$ of a finite module over a commutative ring is
invertible without a noetherian assumption: for generators $x_i$
write $x_i=u(\sum_j a_{ij}x_j)$ and apply the determinant identity
to $I-(a_{ij})u$.  It gives a polynomial identity
$1+u q(u)=0$ on the module, hence the inverse $-q(u)$.
Thus $1-gh$ is a unit of $E$.

Every idempotent of $R$ lifts to $E$.  To see this explicitly,
choose a minimal finite set of generators of $N$ and the resulting
surjection $\pi:A^r\to N$, whose reduction modulo $\mathfrak m$
is an isomorphism.  For $h\in E$ inducing an idempotent on
$N/\mathfrak mN$, choose $H\in M_r(A)$ with $\pi H=h\pi$.
Then $\overline H$ is idempotent, and its characteristic polynomial
is $X^a(X-1)^b$ over $k$.  Henselian factorization gives relatively
prime monic factors $g_0,g_1$ of $\chi_H$ reducing respectively
to $X^a,(X-1)^b$; see
\cite[Lemma 10.153.3(3)--(4)]{StacksHenselFactorSolution}.
They generate the unit ideal in $A[X]$: their resultant is a unit
since its residue is nonzero.  Choose $q\equiv0\pmod{g_0}$ and
$q\equiv1\pmod{g_1}$.  Cayley--Hamilton implies
$q(H)^2=q(H)$.  Since $H$ preserves $\ker\pi$, this descends to
the idempotent $q(h)\in E$, with the required reduction.  The
cases $a=0$ or $b=0$ give the constant lift directly.

A finite-dimensional algebra has a finite decomposition of its
identity into primitive orthogonal idempotents, whose corner
algebras are local.  Lift these idempotents successively, applying
the preceding argument in the complementary summand at each step.
This gives $N=\bigoplus N_i$ with local endomorphism rings:
each corner maps onto its local corner of $R$, with kernel in the
Jacobson radical.  An element invertible in the quotient lifts to
a unit because its two candidate inverse products differ from
$1$ by elements of that radical.  This proves the asserted
decomposition of $M$ over $A$.

Each $M_i$ is finitely presented over $A$.  Finite presentation
and flatness give
\[
 \operatorname{Hom}_A(M_i,M_i)\otimes_A B
     \simeq\operatorname{Hom}_B(M_i\otimes_A B,M_i\otimes_A B).
\]
This follows by applying $\operatorname{Hom}(-,M_i)$ to a finite
presentation and tensoring its kernel diagram with the flat module
$B$.  Taking the kernel of the commutator with $T$ consequently
gives
\[
 E_i\otimes_A B\simeq
 \operatorname{End}_{B[T]}(M_i\otimes_A B),\qquad
 E_i=\operatorname{End}_{A[T]}(M_i).                 \tag{A.s1}
\]
No finite-generation assertion about $E_i$ is used.

Let $R_i$ be the image of $E_i$ on $M_i/\mathfrak mM_i$.
It is a finite-dimensional local $k$-algebra.  Its residue division
algebra $D_i$ is a finite purely inseparable field extension of $k$.
Indeed its center is a finite extension of the separably closed
field $k$, hence purely inseparable and again separably closed.
A central division algebra over that center has a finite separable
splitting field, so is the center itself; this is exactly
\cite[Proposition 11.8.5 and Lemma 11.8.6(4)]{StacksDivisionSolution}.

The image of the ring in (A.s1) on
$(M_i\otimes_A B)/\mathfrak n(M_i\otimes_A B)$ is
$R_i\otimes_k K$.  This equality follows from right exactness
of tensor product and from the injectivity of
$R_i\otimes_k K\to
\operatorname{End}_K((M_i/\mathfrak mM_i)\otimes_k K)$.
The kernel lies in the Jacobson radical by the same finite-module
Nakayama argument as above, now over $B$.
The algebra $R_i\otimes_k K$ is local: its nilpotent ideal
$\operatorname{Jac}(R_i)\otimes_k K$ has quotient
$D_i\otimes_k K$, which is local Artinian because $D_i/k$ is
finite purely inseparable.  For the last claim, a finite set of
purely inseparable generators has a unique possible image in an
algebraic closure of $K$, so the tensor algebra has exactly one
prime ideal.  Thus the ring in (A.s1) is local as well.
\end{proof}

\begin{theorem}[An individual separability tensor over a strict local base]
\label{A:fpqc-local-integral}
\status{Proved}
Let $A$ be strictly henselian local, let $A\to B$ be a flat local
homomorphism, and let $b\in B$ be integral over $A$.  If
\[
 e\in B\otimes_A B,\qquad
 (b\otimes1-1\otimes b)e=0,\qquad \mu(e)=1,
\]
then $b$ is the image of a unique element of $A$.
\end{theorem}
\begin{proof}
Choose a monic equation for $b$.  Lemma~\ref{A:fpqc-finite-module}
places the rank-one $B$-module $B_b$, with $T$ acting by $b$, as a
direct summand of $M_B$ for a finitely presented $A$-module $M$
with $T$ action.  Decompose $M=\bigoplus_i M_i$ by
Proposition~\ref{A:fpqc-indecomposable-basechange}.
Write the components of the inclusion and retraction as
$s_i:B_b\to(M_i)_B$ and $r_i:(M_i)_B\to B_b$.  Then
\[
                 \sum_i r_i s_i=1\quad\hbox{in}\quad
                 \operatorname{End}_{B[T]}(B_b)=B.
\]
Since $B$ is local, one $r_i s_i$ is a unit.  Adjusting that
retraction shows that $B_b$ is a nonzero direct summand of
$(M_i)_B$.  Its endomorphism ring is local, so its only idempotents
are $0$ and $1$; the corresponding nonzero projection must be $1$.
Hence $(M_i)_B\simeq B_b$.

The local flat map $A\to B$ is faithfully flat.  The preceding
isomorphism shows
$\dim_k(M_i/\mathfrak mM_i)=1$, after extending to the residue
field of $B$.  Lift a basis vector; Nakayama gives a surjection
$A\to M_i$.  After tensoring with $B$ this becomes a surjection
between free rank-one $B$-modules, and is therefore an isomorphism.
Flatness and faithful flatness kill its kernel, so $M_i\simeq A$.
On this module $T$ acts by some $a\in A$.  Its base change is
isomorphic to $B_b$ respecting $T$, and therefore $b=a$ in $B$.
Faithful flatness also proves uniqueness.
\end{proof}

\begin{lemma}[The henselization kernel and faithful flatness]
\label{A:fpqc-hensel-kernel}
\status{Proved}
Let $t\in R$ and write $R^h=R^h_{(t)}$.
An element $z\in R$ maps to zero in $R^h$ if and only if
$(1+tc)z=0$ for some $c\in R$.
If $t\in\operatorname{Jac}(R)$ and $R\to D$ is flat with
$R/tR\to D/tD$ faithfully flat, then $R\to D$ is faithfully flat.
Finally, a henselian local ring is henselian along every ideal
contained in its maximal ideal.
\end{lemma}
\begin{proof}
For the second assertion, the image of $\Spec D$ contains every
maximal ideal of $R$, since these contain $t$ and the closed-fiber
map is surjective.  The image of a flat morphism is stable under
generization, so it contains every prime of $R$.

For the first assertion localize $R$ at the multiplicative set
$S=1+tR$.  Then $t\in\operatorname{Jac}(S^{-1}R)$, because
$1+t(c/s)=(s+tc)/s$ is a unit for every $c/s\in S^{-1}R$.
Every element of $S$ is invertible in $R^h$.  The induced map
$S^{-1}R\to R^h$ is flat and is an isomorphism modulo $t$,
by the ind-\'etale construction of henselization
\cite[Lemmas 15.12.1--15.12.2]{StacksHenselPairs}.
It is faithfully flat by the second assertion, so its kernel is
zero.  The kernel of $R\to R^h$ is therefore exactly the kernel
of localization at $S$, giving the stated criterion.

For the last assertion let $(A,\mathfrak m)$ be henselian local
and $I\subset\mathfrak m$.  Every finite $A$-algebra is a finite
product of local rings by
\cite[Lemma 10.153.3(10)]{StacksHenselFactorSolution}.
In each local factor the extension of $I$ lies in its maximal
ideal, and reduction by $I$ preserves that factor's two
idempotents.  Thus reduction induces a bijection on idempotents
in every finite $A$-algebra.  The finite-algebra idempotent
criterion for a henselian pair gives the result
\cite[Lemma 15.11.6]{StacksHenselPairCriterionSolution}.
\end{proof}

\begin{theorem}[Full fpqc descent for henselized $\alpha_p$]
\label{A:hensel-fpqc}
\status{Proved}
For every prime $p$, the functor
\[
 P(R)=\{x\in R^h_{(t)}:x^p=0\}
\]
on $\F_p[t]$-algebras is an fpqc sheaf.  In particular, for every
faithfully flat homomorphism $R\to S$, with no finiteness
hypothesis, the diagram
\[
 P(R)\longrightarrow P(S)\rightrightarrows P(S\otimes_R S)
\]
is an equalizer.
\end{theorem}
\begin{proof}
Proposition~\ref{A:hensel-remaining-gate} proves that $P$ is an
\'etale sheaf, commutes with filtered ring colimits, and is
separated for faithfully flat maps.  We prove existence of descent.

First take a strictly henselian local base $A$, a faithfully flat
map $A\to R$, and a descent datum $b\in P(R)$.  If $t$ is a
unit of $A$, all the relevant henselizations are zero and the
assertion is immediate.  Otherwise $t$ lies in its maximal ideal,
so $(A,tA)$ is henselian by
Lemma~\ref{A:fpqc-hensel-kernel}.  Put $B=R^h_{(t)}$.
The map $A\to B$ is flat and faithfully flat modulo $t$, hence
faithfully flat by that lemma.  The element $b\in B$ satisfies
$b^p=0$.  Functoriality takes the original overlap equality into
$(B\otimes_A B)^h_{(t)}$.  Applying the kernel criterion to
$B\otimes_A B$ gives
\[
             (1+tc)(b\otimes1-1\otimes b)=0
             \quad(c\in B\otimes_A B).             \tag{A.s2}
\]
The multiplication image $u=1+t\mu(c)$ is a unit of $B$, since
$t\in\operatorname{Jac}(B)$.  Thus
$e=(u^{-1}\otimes1)(1+tc)$ has $\mu(e)=1$ and annihilates
$b\otimes1-1\otimes b$.

Choose a prime $\mathfrak q$ of $B$ over the maximal ideal of
$A$.  The homomorphism $A\to B_{\mathfrak q}$ is local and flat.
The image of $e$ in $B_{\mathfrak q}\otimes_A B_{\mathfrak q}$
still has both required properties.  Theorem~\ref{A:fpqc-local-integral}
therefore gives $a\in A$ such that $b=a$ in $B_{\mathfrak q}$.
This equality propagates to all of $B$, as follows.  Set
\[
                D=(B\otimes_A B_{\mathfrak q})^h_{(t)}.
\]
The map $B\to D$ is flat.  Modulo $t$ it is the base change
of the faithfully flat map $A/tA\to B_{\mathfrak q}/tB_{\mathfrak q}$,
and so it is faithfully flat by
Lemma~\ref{A:fpqc-hensel-kernel}, because $t\in\operatorname{Jac}(B)$.
The given cocycle, mapped into $D$, identifies $b$ from its first
factor with $b=a$ from its second factor.  Injectivity of $B\to D$
gives $b=a$ in $B$.  Faithful flatness of $A\to B$ then implies
$a^p=0$.  Hence the datum descends to $P(A)=\alpha_p(A)$.
Notice that no henselian assertion about $B_{\mathfrak q}$ is
needed.

Now let $R\to S$ be arbitrary faithfully flat and let
$\xi\in P(S)$ have equal images on the overlap.  For each prime
$\mathfrak p$ of $R$, choose a geometric point over it and its
strict henselization $A=R_{\mathfrak p}^{\mathrm{sh}}$.
Base change gives a faithfully flat cover $A\to S\otimes_R A$.
The strictly local case supplies $x_{\mathfrak p}\in P(A)$
whose image is the pullback of $\xi$.
The ring $A$ is the filtered colimit of pointed finitely presented
\'etale neighborhoods $R_i$ over $R$.  Finitarity of $P$ first
descends $x_{\mathfrak p}$ to a section $x_i\in P(R_i)$ and,
after enlarging the same index, descends the equality with $\xi$
in $P(S\otimes_R A)$ to equality in $P(S\otimes_R R_i)$.
The images of these \'etale neighborhoods cover $\Spec R$.
On any affine overlap their two sections become equal after the
faithfully flat base change by $S$, and therefore are equal by
fpqc separatedness.  The \'etale sheaf property glues them to
$x\in P(R)$.  Its image equals $\xi$, as may be checked after
the same \'etale cover of $S$.  Uniqueness is separatedness.

This proves all affine faithfully flat equalizers.  Together
with Zariski gluing it is the fpqc sheaf criterion on schemes:
over an affine target a quasi-compact flat surjective morphism
can be refined by a finite union of affine opens in its source,
which is again an affine faithfully flat cover.
\end{proof}

\begin{theorem}[Counterexample to the literal main criterion]
\label{A:main-counterexample}
\status{Proved}
For every prime $p$, over $S=\Spec\F_p[t]$, the functor
$P(R)=\alpha_p(R^h_{(t)})$, extended to schemes by Zariski gluing,
satisfies conditions 1--8 of cote 131, scan 14, with condition 4
on complete noetherian local rings, actual local modularity in 7,
and the locally finite-type interpretation of the tests in 8.
It is not represented by a separated locally quasi-finite
$S$-scheme.  Thus the main representability criterion with this
literal condition 7 is false.
\end{theorem}
\begin{proof}
Proposition~\ref{A:hensel-remaining-gate} verifies every listed
condition except arbitrary fpqc descent, and proves the asserted
nonrepresentability.  Theorem~\ref{A:hensel-fpqc} supplies exactly
that missing condition, so all the hypotheses and the failure of
the conclusion hold simultaneously.

For clarity, the nonrepresentability obstruction occurs already
after base change to $\F_p[[t]]$.  The formal closed-point model
is $\alpha_p$, but a local square-zero test has a section
$\epsilon f$ with $f^p-f=tu$ and $f\equiv0\pmod t$ in its
$t$-henselization, while $f$ is not in the test ring.  The pole
of order one of $tu$ in $\F_p((t))(u)$ rules out an
Artin--Schreier root there.  As proved in the cited proposition,
this excludes an actual finite local modular model and hence a
separated locally quasi-finite representation.  It does not
violate the implication in literal condition 7: on these bases
that implication has no actual finite modular pair as premise.

Proposition~\ref{A:hensel-adic-failure} still applies: this $P$
fails the stronger condition 4 for all noetherian adically complete
rings.  The separate counterexample to L2 is
Theorem~\ref{A:L2-counterexample}, and the criterion with the
explicitly strengthened pro-modularity premise remains
Theorem~\ref{A:pro7-representability}.
\end{proof}
